\section{Введение}

\subsection{Актуальность и мотивация исследования}

Графические ускорители класса NVIDIA GPU заняли центральное место
в современных высокопроизводительных вычислениях~\cite{Gergel2010, GergelStrongin2003}: от обучения крупных
языковых моделей до физического моделирования и научных расчётов.
Программирование GPU на платформе CUDA строится на модели SIMT
(Single Instruction, Multiple Threads), в которой группы из 32~потоков~---
варпы~--- исполняют одну и ту же инструкцию над различными данными~\cite{Lindholm2008, Nickolls2008}.
Ключевым промежуточным представлением в данной экосистеме служит PTX
(Parallel Thread Execution)~--- аппаратно-независимый ассемблер NVIDIA~\cite{PTXISA},
в который компилируется CUDA~C++ и из которого генерируется машинный
код SASS для конкретного поколения GPU.

Несмотря на широкое распространение CUDA, все официальные инструменты
отладки и профилирования~--- Nsight Compute, Nsight Systems,
Compute Sanitizer~--- требуют наличия физического GPU. Это делает их
недоступными в облачных CI-средах, виртуальных машинах и академических
лабораториях с ограниченным бюджетом.

Программный эмулятор PTX~ISA снимает это ограничение: CUDA-ядра
исполняются на стандартном CPU, а встроенный профилировщик собирает
метрики с точностью до отдельной инструкции. Это
определяет следующие области применения:

\begin{enumerate}
  \item \textbf{Среды непрерывной интеграции без GPU.} Автоматическая
        проверка корректности CUDA-ядер при каждом коммите без
        специализированного оборудования.
  \item \textbf{Академические исследования.} Изучение архитектурных
        свойств GPU-программ без доступа к аппаратным счётчикам.
  \item \textbf{Разработка и отладка алгоритмов.} Детерминированное
        пошаговое исполнение с полным контролем над состоянием варпа.
  \item \textbf{Образовательные задачи.} Наглядное представление
        механизмов дивергенции, слияния транзакций памяти и конфликтов
        банков для изучения принципов GPU-архитектуры.
\end{enumerate}


\subsection{Анализ существующих решений}

На сегодняшний день существует ряд инструментов, частично решающих
задачу функциональной эмуляции CUDA-ядер.

\textbf{Ocelot}~\cite{Diamos2010}~--- один из первых функциональных
PTX-интерпретаторов с поддержкой перехвата через CUDA~Runtime~API.
Его ключевое ограничение~--- привязка к PTX~ISA версий 1.x--2.x (CUDA~4.x)
и отсутствие поддержки директив PTX~7.x+, в частности инструкций
\texttt{ldmatrix}, \texttt{wmma} и атрибута \texttt{.address\_size~64}.
Проект не поддерживается с 2013~года.

\textbf{GPGPU-Sim}~\cite{Bakhoda2009}~--- timing-accurate микроархитектурный
симулятор. Поддерживает PTX~ISA только до версии~4.x и ориентирован
на оценку латентностей, а не функциональную корректность.

\textbf{Compute Sanitizer}~--- официальный инструмент NVIDIA
для обнаружения гонок данных, ошибок памяти и нарушений барьерной
синхронизации. Требует наличия физического GPU и не предоставляет
метрик уровня инструкций PTX.

Актуального PTX-эмулятора, поддерживающего ISA~7.x+, с интегрированным профилировщиком и прозрачным перехватом CUDA~Runtime~API без модификации бинарного файла, не существует~--- этот пробел и заполняет данная работа.

\subsection{Цели и задачи}

\textbf{Цель работы}~--- разработка программного эмулятора PTX~ISA,
обеспечивающего функционально корректное исполнение CUDA-ядер без GPU
с возможностью сбора детальных метрик производительности на уровне
инструкций.

Выбор PTX~ISA в качестве уровня эмуляции обусловлен тем, что PTX
является стабильным публичным контрактом NVIDIA~\cite{PTXISA}: компилятор гарантирует
его синтаксическую и семантическую совместимость между поколениями GPU,
тогда как SASS~--- бинарный машинный код~--- документируется значительно хуже
и изменяется с каждым поколением архитектуры.

Для достижения указанной цели поставлены следующие \textbf{задачи}:
\begin{enumerate}
  \item Реализовать механизм прозрачного перехвата CUDA~Runtime~API
        через \texttt{LD\_PRELOAD} с извлечением PTX посредством \texttt{cuobjdump}.
  \item Разработать однопроходный инкрементный синтаксический анализатор
        PTX-текста, поддерживающий подмножество PTX~ISA~7.x.
  \item Построить таблицу диспетчеризации и исполнители инструкций
        с формальным описанием поведения инструкций.
  \item Реализовать алгоритм PDOM-реконвергенции для управления
        дивергентным потоком варпа.
  \item Разработать систему метрик профилировщика с автогенерацией
        исполнителей из JSON файла.
  \item Реализовать инструмент генерации самодостаточного HTML-отчёта
        с аннотацией строк исходного кода.
  \item Верифицировать корректность эмулятора на семи интеграционных
        тестах с реальными CUDA-ядрами.
\end{enumerate}


\subsection{Научная новизна}

Научная новизна работы заключается в следующем.

\begin{enumerate}
  \item Впервые реализован прозрачный механизм перехвата CUDA~Runtime~API
        посредством \texttt{LD\_PRELOAD} без модификации целевого бинарного
        файла с поддержкой PTX~ISA~7.x; дано формальное обоснование порядка
        разрешения символов динамическим загрузчиком.
  \item Разработан и верифицирован однопроходный синтаксический анализатор
        подмножества PTX~ISA~7.x, корректность которого подтверждена
        на семи CUDA-ядрах и тридцати одном модульном тесте.
  \item Описано поведение инструкций подмножества PTX~ISA в виде
        формальных правил, позволяющих верифицировать корректность
        исполнителей без обращения к аппаратному обеспечению.
  \item Предложена архитектура расширяемой системы сбора метрик
        с автогенерацией исполнителей из JSON-файла, при которой
        добавление новой метрики не требует модификации кода на C++.
  \item Получена формальная модель коэффициента слияния обращений к
        глобальной памяти \(\kappa = T_{\mathrm{ideal}} / T_{\mathrm{actual}}\)
        с учётом размера элемента доступа.
\end{enumerate}


\subsection{Практическая значимость}

Разработанный инструмент имеет непосредственное практическое применение:

\begin{itemize}
  \item \textbf{Внедрение в конвейеры непрерывной интеграции (CI).}
        Эмулятор запускается как обычный Linux-процесс; для включения
        в \texttt{GitHub Actions} или \texttt{Jenkins} достаточно
        однострочного скрипта \texttt{cuemu ./binary}.
  \item \textbf{Детальное профилирование.} HTML-отчёт с аннотацией
        строк исходного кода позволяет определить «горячие» инструкции
        и оптимизировать шаблоны обращения к памяти без запуска на GPU~\cite{Meerov2009}.
  \item \textbf{Образовательная ценность.} Пошаговая трассировка
        маски дивергенции обеспечивает наглядную демонстрацию
        механизма SIMT-исполнения.
\end{itemize}

Инструмент не требует специализированного оборудования: PTX извлекается утилитой \texttt{cuobjdump} из набора NVIDIA, ядра исполняются полностью на CPU.

\newpage
